% =====================================================================
% execucao-poa.tex -- registro de execucao
%
% Primeiro experimento da familia "avaliacao de recurso": medir o erro
% do NASA POWER contra a irradiancia no plano do arranjo que o ONS
% publica em `restricao_coff_fotovoltaica_detail`.
%
% E um registro de EXECUCAO, nao um capitulo de resultado. Guarda o que
% deu errado no caminho e como se descobriu, porque tres dos quatro
% achados sao defeitos de dado que nenhuma leitura do codigo pegaria --
% so apareceram ao medir.
%
% Nao entra no `reporte.tex`. Compila sozinho, e em lualatex, nao em
% pdflatex: a sem serifa aqui e a Urbanist, que vive como TTF em
% `fontes/urbanist/` e so `fontspec` alcanca. Rode de dentro de
% `report/`, que e onde o `Path` do `fontspec` esta ancorado:
%   latexmk -lualatex execucao-poa.tex
% Em pdflatex o documento ainda compila -- a folha de estilo cai de volta
% na IBM Plex Sans -- mas nao e a saida pretendida.
% =====================================================================

% A folha de estilo pressupoe `book`: usa `\chaptertitlename` e os
% comandos de capitulo do `tocloft`. Em `article` ela nao carrega.
\documentclass[10pt,a4paper,oneside,twocolumn]{book}

\usepackage{iftex}
\ifPDFTeX\usepackage[utf8]{inputenc}\fi
\usepackage[brazilian]{babel}
% Mancha e medianiz do `reporte.tex`: 17,2 cm de mancha em duas colunas
% de 8,25 cm. Duas colunas resolvem a exigencia de aproveitar a folha sem
% alongar a linha -- em coluna unica a mesma mancha daria 73 caracteres.
\usepackage[a4paper,top=2.0cm,bottom=2.2cm,left=1.9cm,right=1.9cm,
            columnsep=0.7cm]{geometry}
\usepackage{estilo-reporte}

% A pagina termina onde o texto termina.
\raggedbottom
\pagestyle{plain}
\setlength{\parskip}{0.4em}
\setlength{\parindent}{0pt}

\autoria{João Leonardi da Silva Melo}{joao\_leonardi.melo@somosicev.com}

\begin{document}

\abertura{Erro do NASA POWER contra a POA medida}
         {Registro de execução · 5 de setembro de 2026 ·
          \texttt{notebooks/br/\_poa\_erro.py}}

\section*{O que se quis medir}

A etapa~0 estabeleceu que \texttt{val\_irradianciaverificado} é medida no plano
do arranjo, e não na horizontal, e mediu que o NASA POWER fica abaixo dela.
Faltava fechar a conta: transpor o GHI do POWER para o plano e ver o que sobra.

O interesse não é acadêmico. Trocar NASA POWER por NSRDB está listado como
trabalho seguinte, e o argumento escrito é sobre resolução --- a célula passa de
cerca de 50\,km para 2\,km. Resolução não é erro. Sem medir o erro atual não há
como dizer se a troca compensa.

\section*{Desenho}

Quatro geometrias sobre a mesma série, e a diferença entre elas é o resultado.

\begin{ficha}
  \item[0°]          controle: a transposição tem de devolver o próprio GHI.
  \item[$|$lat$|$]   a regra a priori para arranjo fixo, sem olhar o medido.
  \item[Ajustada]    a inclinação fixa que minimiza o erro da própria usina.
  \item[Rastreador]  eixo único horizontal norte-sul, faixa de 60°, com
                     \emph{backtracking} e GCR de 0,35.
\end{ficha}

O erro em $|$lat$|$ é o que um usuário do POWER teria sem verdade de campo. O
erro na ajustada é o piso do arranjo fixo. A distância entre os dois é o custo
de não conhecer a inclinação. O rastreador entrou depois, como hipótese que a
própria distância obrigou a testar.

\section*{Dado}

Janela de 2025, que é o que o cache horário cobre. O pareamento é interno ao
banco: \texttt{br.pv\_detail} casa com \texttt{br.plant} por \texttt{ceg\_core}
em 560 de 560 usinas, e 558 têm geometria. Sem casamento de nome e sem cadastro
externo.

\begin{ficha}
  \item[QC noturno]  554 de 560 usinas com mediana de irradiância $\le$ 5 W/m²
                     entre 22\,h e 3\,h.
  \item[Dia válido]  ao menos 44 dos 48 intervalos, e mais de 1 kWh/m².
  \item[Usina]       ao menos 60 dias válidos no ano.
  \item[Resultado]   \num{69550} dias-usina em 338 usinas; \num{66917} dias em
                     237 usinas depois do corte de 60 dias.
  \item[Células]     66 pontos do POWER, todos presentes no cache.
\end{ficha}

\notas{Os 101 descartes por cobertura não são falha de método: 66 dessas usinas
têm 30 dias ou menos de série limpa em 2025. Entre as 237 retidas, a mediana é
de 321 dias e 183 passam de 240.}

\section*{Quatro defeitos, todos encontrados medindo}

\campo{Cache com fuso misturado} A primeira rodada deu 29\% de subestimação,
contra os 14,8\% que a etapa~0 tinha medido \emph{sem} transposição. Transpor
para o plano não pode piorar. Rastreando: nove das 90 células do cache têm
\texttt{time\_standard=\textquotesingle LST\textquotesingle} no cabeçalho, hora
solar local, e todo leitor do diretório assume UTC. Nessas células até 14,2\% do
GHI anual caía em intervalo com o sol abaixo do horizonte.

A causa é uma linha: \texttt{\_power\_h.py} não passa
\texttt{\&time-standard=UTC}, e a API do POWER responde em LST por omissão. O
\texttt{\_power\_h\_todas.py}, que faz certo, pulou os nove arquivos porque
testava apenas \texttt{os.path.exists}. Os dois foram corrigidos, e o segundo
agora valida o cabeçalho em vez da existência.

\campo{Componentes que não fecham} \texttt{DNI$\cdot$cos(z) + DHI} do POWER não
reproduz o GHI publicado. Entregá-las ao pvlib perde a diferença antes de
qualquer inclinação: a 0°, onde a transposição deveria devolver o GHI, saía 8\%
baixa. Reescalar as componentes para fechar não resolve --- o fator satura no
teto em 27,7\% dos intervalos com sol. A saída foi descartar DNI e DHI
publicados e decompor o GHI por Erbs, que fecha por construção. O fechamento das
componentes publicadas ficou registrado como diagnóstico: mediana 0,9565, com
p10 de 0,9547 e p90 de 0,9590.

\notas{A dispersão que se via antes --- fechamento de 0,75 a 0,95 conforme a
célula --- era o defeito de fuso. Corrigido o fuso, a falta de
fechamento é uniforme, e é característica do produto.}

\campo{Carimbo da hora} O valor horário do POWER é média da hora, e a posição
solar precisa ser calculada no meio dela. Calculada no instante do carimbo, o
\texttt{airmass} sai NaN em 29,1\% dos intervalos com sol.

\campo{Deslocamento noturno do piranômetro, e um controle na janela errada}
Achado depois, ao investigar sete usinas de forma diurna anômala. O controle de
qualidade rodava a mediana noturna sobre o período inteiro do banco, e a análise
é de 2025: Santa Luzia~1 e~V passavam com 2,8 W/m², porque o sensor estava são
em 2024 e 2026 e deslocado em 323,8 durante todo o~2025 --- cerca de
3,9 kWh/m²/dia de artefato somados ao denominador. O controle agora corre na
mesma janela da análise.

E a exclusão por limiar foi trocada por subtração. A média noturna da frota é
contínua --- mediana 4,6 W/m², p75 15,5, p90 69,1 --- e não tem vale onde
cortar; qualquer limiar custaria de um quarto a metade das usinas. Subtrai-se a
mediana noturna da própria usina naquele mês, o que cobre o caso de sensor que
degrada no meio do ano, e exclui-se apenas quando o resíduo depois da subtração passa de
20 W/m², que é falha intermitente e não deslocamento.

\conclusao{Nenhum dos quatro apareceria em revisão de código. Os quatro
apareceram porque um número divergiu de outro já estabelecido. Dois são do
dado --- o cache em fuso local e as componentes que não fecham --- e dois são
de tratamento: a posição solar no carimbo e a janela do controle.}

\section*{Resultado}

Controle aprovado nas 237 usinas: \texttt{poa(0°)/GHI} entre 0,9999 e 1,0000.

\vspace{0.4em}
\tabfont
\begin{tabular}{@{}>{\rag\arraybackslash}p{3.3cm}rrr@{}}
\toprule
\cab{Modelo} & \cab{Viés} & \cab{MAE} & \cab{Razão} \\
\midrule
GHI horizontal            & $-0{,}831$ & 1,228 & 1,229 \\
Fixo em $|$lat$|$         & $-0{,}660$ & 1,150 & 1,159 \\
Fixo no melhor ângulo     & $-0{,}721$ & 1,126 & 1,142 \\
Rastreador de eixo único  & $+0{,}459$ & 0,998 & 1,000 \\
\addlinespace[.4ex]
Escolha por usina         & ---        & \textbf{0,705} & 1,014 \\
\bottomrule
\end{tabular}
\notas{Viés e MAE em kWh/m²/dia; razão é medido sobre modelo, mediana entre
usinas. \num{237} usinas, \num{66917} dias-usina, 2025.}

\vspace{0.6em}
A pergunta era o erro do POWER, e a resposta é que quase todo ele era geometria
não modelada. A transposição para plano fixo cobre um terço da lacuna; o
rastreador cobre o resto. O viés vai de $-0{,}83$ a $+0{,}46$ kWh/m²/dia, e
escolher a geometria por usina derruba o MAE em 43\%.

\section*{Duas validações}

\campo{Coerência por conjunto} Usinas do mesmo conjunto compartilham tecnologia.
Em 28 conjuntos com duas ou mais usinas medidas, cobrindo 229 usinas, 86\%
recebem rótulo unânime --- contra 13\% esperados se o rótulo fosse sorteio com a
mesma proporção. A margem também é larga: mediana de 110\% de diferença de MAE
entre as duas geometrias, e apenas 2\% das usinas ficam em empate abaixo de 5\%.

\campo{Janaúba} A etapa~0 usou Janaúba como controle interno: lê razão 1,00
enquanto as outras leem cerca de 1,24, e concluiu diferença de montagem do
sensor. Chegando por outro caminho --- ajuste da forma diurna, e não da razão
anual --- a classificação põe 18 das 27 usinas de Janaúba como fixas, invertendo
a proporção da frota. As razões remedidas são 1,023 em Janaúba contra 1,240 nas
demais, reproduzindo a etapa~0.

\section*{O que não se pode afirmar}

\begin{objecao}
Que 66\% da frota usa rastreador. A classificação separa as duas classes com
margem mediana de 110\% em MAE e é coerente dentro do conjunto, mas não
acompanha a safra de entrada em operação: a fração classificada
como rastreador por ano vai 1,00 · 1,00 · 0,91 · 0,38 · 0,83 · 0,45 · 0,93, sem
tendência, e 100\% de rastreador em 2017 é implausível. O porte mediano mal
difere entre os grupos, 39,5 contra 33,5 MW. A geometria de rastreador captura
uma propriedade da usina e por usina --- a coerência de conjunto sustenta isso ---, mas o rótulo
pode estar absorvendo outras diferenças de sítio junto com a montagem. Lê-lo como
inventário de tecnologia exige confronto com cadastro de tipo de montagem, que
não se procurou ainda.
\end{objecao}

\section*{Remedição das saídas da etapa 0}

\texttt{\_transpor.py}, \texttt{\_otimo.py} e \texttt{\_tilt\_pareado.py} leem o
mesmo cache assumindo UTC, e sustentam as afirmações sobre inclinação ajustada e
ótima. A amostra de 25 usinas do \texttt{\_power\_h.py} é exatamente a que gerou
as nove células deslocadas, então a contaminação ali era alta. Os três foram
refeitos, com as saídas anteriores guardadas em
\texttt{data/processed/\_antes\_correcao\_lst/}.

\vspace{0.4em}
\tabfont
\begin{tabular}{@{}>{\rag\arraybackslash}p{3.6cm}rr@{}}
\toprule
\cab{Grandeza} & \cab{Antes} & \cab{Depois} \\
\midrule
Inclinação ajustada, 25 usinas   & 20,0°  & 17,5° \\
Ajustado $-$ ótimo, 25 usinas    & $+5{,}0$° & $+0{,}0$° \\
Ajustado $-$ ótimo, 174 pareadas & $+2{,}5$° & $+0{,}0$° \\
Em ou abaixo do ótimo            & 45\%   & 53\% \\
Janaúba, $p$ de Mann-Whitney     & 0,015  & 0,056 \\
\bottomrule
\end{tabular}
\notas{Inclinação ajustada mudou em 24 das 25 usinas de conveniência e em 49
das 174 pareadas. O desenho pareado absorveu melhor o defeito, como se espera
dele.}

\vspace{0.6em}
Duas conclusões mudam. A de que os projetos reais inclinam acima do ótimo
calculado não se sustenta: a mediana da diferença passa a ser exatamente zero, pelos dois
caminhos. E a previsão sobre Janaúba deixa de ser significativa ao limiar
convencional --- as medianas não mudam, mas o $p$ cruza 0,05.

\campo{Cadastro de montagem} Procurado e não encontrado. ANEEL SIGA, MMGD e
\texttt{mmgd\_pv\_tec}, RALIE de usina e de unidade geradora, ONS
\texttt{coff\_fv\_detail}, \texttt{coff\_fotovoltaica} e
\texttt{modalidade\_usina}: nenhum traz o tipo de montagem. O
\texttt{mmgd\_pv\_tec} chega mais perto, com área de arranjo e contagem de
módulos, mas cobre micro e minigeração, não a frota centralizada. A classificação
fixo contra rastreador segue sem validação externa.

\section*{Artefatos}

\begin{ficha}
  \item[Código]   \texttt{notebooks/br/\_poa\_erro.py}; correções em
                  \texttt{\_power\_h.py} e \texttt{\_power\_h\_todas.py}.
  \item[Saídas]   \texttt{data/processed/poa\_erro\_por\_usina.csv}
                  (237 × 28) e
                  \texttt{poa\_erro\_pareado.parquet} (\num{66917} × 12).
  \item[Dado]     nove células rebaixadas em UTC; as originais em LST estão em
                  \texttt{data/raw/power\_hourly\_lst\_backup/}.
\end{ficha}

\end{document}
